\documentclass[11pt]{article}
\usepackage[margin=1in]{geometry}
\usepackage[utf8]{inputenc}
\usepackage[T1]{fontenc}
\usepackage{lmodern}
\usepackage{amsmath,amssymb,amsthm}
\DeclareUnicodeCharacter{2200}{\ensuremath{\forall}}
\DeclareUnicodeCharacter{2192}{\ensuremath{\to}}
\DeclareUnicodeCharacter{2243}{\ensuremath{\simeq}}
\DeclareUnicodeCharacter{2097}{\ensuremath{_{l}}}
\DeclareUnicodeCharacter{2045}{\ensuremath{[\![}}
\DeclareUnicodeCharacter{2046}{\ensuremath{]\!]}}
\DeclareUnicodeCharacter{22A4}{\ensuremath{\top}}
\DeclareUnicodeCharacter{03C1}{\ensuremath{\rho}}
\DeclareUnicodeCharacter{00AC}{\ensuremath{\neg}}
\usepackage[title]{appendix}
\usepackage[colorlinks=true,linkcolor=blue,citecolor=blue,urlcolor=blue,pdfauthor={P. Graham},pdftitle={No post-Lie algebra structures on pairs of perfect and nilpotent Lie algebras}]{hyperref}
\newtheorem{theorem}{Theorem}[section]
\newtheorem{proposition}[theorem]{Proposition}
\newtheorem{corollary}[theorem]{Corollary}
\newtheorem{lemma}[theorem]{Lemma}
\newtheorem{question}[theorem]{Question}
\theoremstyle{definition}
\newtheorem{definition}[theorem]{Definition}
\theoremstyle{remark}
\newtheorem{remark}[theorem]{Remark}
\setlength{\emergencystretch}{3em}
\newcommand{\g}{\mathfrak{g}}
\newcommand{\n}{\mathfrak{n}}
\newcommand{\s}{\mathfrak{s}}
\newcommand{\h}{\mathfrak{h}}
\DeclareMathOperator{\ad}{ad}
\DeclareMathOperator{\tr}{tr}
\DeclareMathOperator{\Der}{Der}
\DeclareMathOperator{\End}{End}
\DeclareMathOperator{\adj}{adj}
\title{No post-Lie algebra structures on pairs of perfect and nilpotent Lie algebras}
\author{P. Graham\\ \small\url{https://orcid.org/0009-0007-8493-1505}}
\date{November 10, 2026}
\begin{document}
\maketitle
\begin{abstract}
Let $\g$ and $\n$ be finite-dimensional Lie algebras on the same vector
space over a field of characteristic zero, with $\n$ nilpotent. We prove
that if the pair $(\g,\n)$ carries a post-Lie algebra structure all of
whose right multiplications have trace zero, then $\g$ is solvable. Hence
no nonzero perfect Lie algebra $\g$ admits a post-Lie algebra structure on
$(\g,\n)$ with $\n$ nilpotent. This answers a question of Burde and Dekimpe
in the negative and proves the nilpotent case of a conjecture of Burde,
Dekimpe and Monadjem. Nonexistence was known for abelian $\n$, by a theorem of
Helmstetter on left-symmetric algebras; for two-step nilpotent $\n$, by
Burde and Dekimpe; and for the complex perfect
non-semisimple $\g$ of dimensions five and six, by Burde and Dekimpe and by
Burde, Dekimpe and Monadjem. The proof extends Helmstetter's argument. A
post-Lie algebra structure embeds $\g$ into the holomorph of $\n$.
Conjugations chosen by Whitehead's first lemma along the lower central
series of $\n$ make the translation part of this embedding vanish on a
Levi factor of $\g$, while its determinant stays equal to one, as in
Kriebel's treatment of semisimple $\n$. The proof is
formalized in Lean with Mathlib, assuming Levi's theorem and Whitehead's
first lemma as stated in Weibel's textbook.

This work has not yet been reviewed by experts.
The proofs, the Lean formalization and the text were produced by AI models
(from OpenAI and Anthropic), working under the author's direction.
\end{abstract}
\medskip
\noindent\textit{Keywords:} post-Lie algebra, perfect Lie algebra,
nilpotent Lie algebra, left-symmetric algebra, NIL-affine action.

\noindent\textit{MSC 2020:} 17B30, 17B60, 17D25.

\section{Introduction}

\subsection{Main results}
Throughout, $k$ is a field of characteristic zero and all Lie algebras
are finite-dimensional over $k$. Following Burde and Dekimpe
\cite[Definition 2.1]{BD}, let $\g=(V,[\,,])$ and $\n=(V,\{\,,\})$ be two
Lie brackets on one vector space $V$. A \emph{post-Lie algebra structure}
on the pair $(\g,\n)$ is a bilinear product $x\cdot y$ on $V$ such that,
for all $x,y,z\in V$,
\begin{align}
x\cdot y-y\cdot x&=[x,y]-\{x,y\},\label{eq:pl1}\\
[x,y]\cdot z&=x\cdot(y\cdot z)-y\cdot(x\cdot z),\label{eq:pl2}\\
x\cdot\{y,z\}&=\{x\cdot y,z\}+\{y,x\cdot z\}.\label{eq:pl3}
\end{align}
We write $L(x)y=x\cdot y$ and $R(x)y=y\cdot x$. By \eqref{eq:pl3} each
$L(x)$ is a derivation of $\n$, and by \eqref{eq:pl2} the map
$L\colon\g\to\Der(\n)$ is a representation of $\g$ \cite[\S2]{BD}. When
$\n$ is abelian, \eqref{eq:pl1} and \eqref{eq:pl2} say that $(V,\cdot)$ is
a left-symmetric algebra (an algebra whose associator
$(x,y,z)=(x\cdot y)\cdot z-x\cdot(y\cdot z)$ is symmetric in $x$ and $y$)
whose commutator is the bracket of $\g$.

\begin{theorem}\label{thm:main}
Let $x\cdot y$ be a post-Lie algebra structure on a pair $(\g,\n)$ with
$\n$ nilpotent. If $\tr R(x)=0$ for every $x\in V$, then $\g$ is
solvable.
\end{theorem}

\begin{corollary}\label{cor:perfect}
Let $\g\neq0$ be perfect, that is, $[\g,\g]=\g$, and let $\n$ be
nilpotent. Then there is no post-Lie algebra structure on $(\g,\n)$.
\end{corollary}

For $\g=0$ the zero product is a post-Lie algebra structure, so the
hypothesis $\g\ne0$ is needed. The trace hypothesis in
Theorem~\ref{thm:main} cannot be dropped: the product of matrices
$x\cdot y=xy$ on $V=\mathfrak{gl}_m(k)$ is a post-Lie algebra structure on
$(\mathfrak{gl}_m(k),\n)$ with $\n$ abelian, since associativity gives
\eqref{eq:pl2}, and $\mathfrak{gl}_m(k)$ is not solvable for $m\ge2$; here
$\tr R(x)=m\tr x$.

\subsection{The question and prior work}
Post-Lie algebras were introduced by Vallette \cite{Va}. As Burde and
Dekimpe recall \cite[\S1]{BD}, post-Lie algebra structures on pairs of Lie
algebras were introduced independently by Burde, Dekimpe and Vercammen
\cite{BDV}, in the study of NIL-affine actions of Lie groups, that is,
affine actions on a nilpotent Lie group $N$ by elements of
$N\rtimes\mathrm{Aut}(N)$.

Burde and Dekimpe \cite{BD} studied their existence in terms of the
structure of $\g$ and $\n$. For perfect $\g$ they proved: ``Proposition
3.3. Let $(\g,\n)$ be a pair of Lie algebras over a field of
characteristic zero, where $\g$ is perfect and $\n$ is abelian or 2-step
nilpotent. Then there is no post-Lie algebra structure on $(\g,\n)$.'' In
its proof the abelian case is taken from Segal \cite[Theorem 1 and
Proposition 1]{Se} and, over $\mathbb{C}$, from Helmstetter
\cite[Proposition 20]{He}, and the two-step case is reduced to it by
\cite[Proposition 4.2]{BDV}. They also recall that $\n$ cannot be solvable
and not nilpotent \cite[Proposition 4.4]{BDV} (their Proposition 3.4), and
that $\g$ cannot be semisimple when $\n$ is solvable \cite[Theorem
4.2]{BD13} (their Theorem 3.2). This led them to ask:

\begin{question}[{\cite[Question 3.5]{BD}}]
Let $(\g,\n)$ be a pair of Lie algebras, where $\g$ is perfect and $\n$ is
nilpotent. Does there exist a post-Lie algebra structure on $(\g,\n)$?
\end{question}

They showed that no structure exists for
$\g=\mathfrak{sl}_2(\mathbb{C})\ltimes V(2)$, the semidirect product with
the two-dimensional irreducible module \cite[Proposition 3.6]{BD}. In a
report of joint work with Dekimpe, Burde then stated two conjectures
\cite[p.~2689]{OWR}: ``Conjecture. Let $(\g,\n)$ be a pair of Lie
algebras, where $\g=\mathfrak{sl}_n(\mathbb{C})\ltimes V(n)$ and $\n$ is
nilpotent. Then there is no post-Lie algebra structure on $(\g,\n)$,''
where $n$ is an integer and $V(n)$ the natural representation of
$\mathfrak{sl}_n(\mathbb{C})$, and, as ``a more general conjecture'':
``Conjecture. Let $(\g,\n)$ be a pair of Lie algebras, where $\g$ is
perfect non-semisimple, and $\n$ is nilpotent. Then there is no post-Lie
algebra structure on $(\g,\n)$.''

Burde, Dekimpe and Monadjem \cite{BDM} studied the pairs with $\g$ perfect
and not semisimple, in seven cases for $\n$, of which case (b) is ``$\n$
is nilpotent non-abelian'' \cite[\S3]{BDM}. In their introduction they
write that for the cases ``where $\g$ is perfect non-semisimple, and $\n$
is either nilpotent, simple, or semisimple, we are not able [to] solve the
existence question in general. We conjecture that there do not exist
post-Lie algebra structures in these cases.'' In case (b) they proved
nonexistence for the two six-dimensional complex perfect Lie algebras that
are not semisimple, $\mathfrak{sl}_2(\mathbb{C})\ltimes V(3)$ and
$\mathfrak{sl}_2(\mathbb{C})\ltimes\mathfrak{n}_3(\mathbb{C})$, where
$\mathfrak{n}_3(\mathbb{C})$ is the three-dimensional Heisenberg algebra
\cite[Proposition 3.2]{BDM}.

Corollary~\ref{cor:perfect} answers Question 3.5 of \cite{BD} in the
negative, and proves case (b) of the conjecture of \cite{BDM} and both
conjectures of \cite[p.~2689]{OWR}, over every field of characteristic
zero.

The case $\n$ abelian of Theorem~\ref{thm:main} is due to Helmstetter. A
left-symmetric algebra is \emph{complete} if every $I+R(x)$ is bijective
\cite[Definition 3.30]{Bu}, and in characteristic zero this holds if and
only if every $R(x)$ is nilpotent, and if and only if $\tr R(x)=0$ for
all $x$ \cite[Theorem 3.31]{Bu}, a theorem of Segal \cite{Se}. Helmstetter
calls a left-symmetric algebra nilpotent when every right multiplication
is nilpotent \cite[Definition 17]{He}, and proved: ``Si $A$ est une
alg\`ebre sym\'etrique \`a gauche nilpotente, l'alg\`ebre de Lie $A$ est
r\'esoluble'' (if $A$ is a nilpotent left-symmetric algebra, its Lie
algebra is solvable) \cite[Proposition 20]{He}. The proof of
Helmstetter's Corollary 21, that for a finite-dimensional left-symmetric algebra ``il
est impossible que l'alg\`ebre de Lie $A$ soit \'egale \`a l'alg\`ebre de
Lie d\'eriv\'ee $[A,A]$'' (the Lie algebra $A$ cannot equal its derived
algebra), deduces nilpotency from $\tr R_a=0$ for every $a$. Fox restates
these results, as due to Helmstetter, in \cite[Lemma 5.3]{Fo}: ``(1) The
underlying Lie algebra of a complete finite-dimensional LSA over $k$ is
solvable.'' Burde recalls the geometric origin: if a Lie group $G$ acts
simply transitively by affine transformations of $\mathbb{R}^n$, then
``Auslander proved that in this case $G$ is solvable'' \cite{Au},
\cite[Remark 2.28]{Bu}.

Several analogues concern groups. A Lie skew brace is a Lie group with a
second Lie group structure on the same manifold, linked by the skew-brace
identity, and its tangent object is a post-Lie algebra \cite[\S1]{DL}.
Damele and Loi proved, in their Theorem 1.1(S2): ``Let $(G,\cdot,\circ)$
be a connected LSB [Lie skew brace]. \dots\ If $(G,\cdot)$ is solvable,
then $(G,\circ)$ is solvable'' \cite{DL}. Of the existence question for
nilpotent $\n$ and perfect $\g$, which they attribute to \cite{BDM}, they
write: ``While we are currently not able to answer this question
definitively, we can nevertheless state that any such PLA [post-Lie
algebra], if it exists, would not be integrable into an LSB due to
condition (S2) in Theorem 1.1'' \cite[\S1]{DL}. Theorem~\ref{thm:main}
needs no Lie group and no integration. For finite groups, Byott proved
``if $G$ is nilpotent then $\Gamma$ is soluble'' for a Hopf--Galois
structure of type $G$ on a Galois extension with group $\Gamma$
\cite[Theorem 1]{By}; for skew braces this reads ``Let $A$ be a finite
skew brace whose additive group $(A,\cdot)$ is nilpotent. Then $(A,\circ)$
is solvable'' \cite[Corollary 1.23]{SV}, and Damele has extended it to
finite skew braces whose additive derived subgroup is nilpotent \cite{Da}.
On simply connected nilpotent Lie groups, Der\'e and Origlia recall that
``if a group acts simply transitively on a 1-connected nilpotent Lie group,
it must be solvable'' \cite[\S1]{DO}.

For semisimple $\n$, Kriebel's preprint gives a proof of nonexistence for
perfect non-semisimple complex $\g$ \cite[Theorem 1 and Corollary 3]{Kr}.
There ``the integrated left-right action has a holomorphic infinitesimal
determinant. Unimodularity makes this determinant invariant''
\cite[abstract]{Kr}, and Kriebel credits Baues's determinant relative
invariant of an \'etale affine representation as its predecessor
\cite[Theorem 5.19 and Corollary 5.20]{Ba}.

\subsection{Proof overview and what is new}
Helmstetter's proof, in the form given by Fox \cite[proof of Lemma
5.3]{Fo}, runs as follows. If $\g$ is not solvable it contains a nonzero
semisimple subalgebra $\s$; the identity map $\s\to V$ is a $1$-cocycle for
the representation $L$, so by Whitehead's first lemma there is $a\in V$
with $x=L(x)a=R(a)x$ for $x\in\s$. Hence $I-R(a)$ is singular, which
contradicts completeness (take $x=-a$).

For nilpotent $\n$ both steps need a replacement. By a theorem of Burde,
Dekimpe and Vercammen \cite[Proposition 2.13]{BDV}, a post-Lie algebra
structure is the same as an embedding $x\mapsto(x,L(x))$ of $\g$ into the
holomorph $\n\rtimes\Der(\n)$. Its first component, the translation part,
is the identity map $P=\mathrm{id}_V$, which is a cocycle for $L$ with
respect to the bracket of $\n$ (Section~\ref{sec:setup}). Conjugating the
embedding by the exponential of $-\ad(z,0)$ for $z\in\n$ changes $P$ to a
new cocycle $P^z$. We show two things.
\begin{itemize}
\item When $\tr L=\tr\ad_\g$, which is equivalent to $\tr R=0$, every
conjugation preserves $\det P$ (Lemma~\ref{lem:det}). This step follows
the determinant method Kriebel used for semisimple $\n$ \cite{Kr}, after
Baues \cite{Ba} and Helmstetter's characteristic polynomial: there the
unimodularity of $\g$ makes a determinant invariant, and here the trace
hypothesis does. What is new is that for nilpotent $\n$ the invariance is
an identity of polynomials, with no integration. The proof puts $tz$ in
place of $z$, so that the translation part becomes a matrix $P_t$ with
entries in $k[t]$, and shows that $\det P_t$ has derivative zero by
Jacobi's formula and a trace identity extracted from the cocycle equation.
For $\n$ abelian, $P_t=I+tR(z)$, and the lemma says that Helmstetter's
characteristic polynomial $\det(I+R(x))$ \cite[\S3]{Fo} equals one.
\item Let $\n=\n^1\supseteq\n^2\supseteq\dots\supseteq\n^{r+1}=0$ be the
lower central series and $\s$ a Levi factor of $\g$. Whitehead's first
lemma, applied to the successive quotients $\n^i/\n^{i+1}$, gives elements
$z_1,\dots,z_r$ whose conjugations move $P(\s)$ into
$\n^2,\dots,\n^{r+1}=0$ (Lemma~\ref{lem:whitehead}). In the abelian case
one application suffices. When $\n$ is not abelian the cocycle equation is
nonlinear and becomes linear only modulo the next term of the series, so
the lemma is applied $r$ times; this step has no abelian analogue.
\end{itemize}
After the $r$ conjugations the determinant is still one, but the cocycle
vanishes on $\s\ne0$, a contradiction.

The proof uses Levi's theorem and Whitehead's first lemma, quoted in
Section~\ref{sec:lean}. Everything else is proved in a Lean formalization
\cite{Lean,Mathlib}, described there; the two theorems enter it as named
axioms.

\section{Cocycles and conjugation}\label{sec:setup}
Let $\n=(V,\{\,,\})$ be a Lie algebra and $\h=\n\rtimes\Der(\n)$ its
holomorph, the vector space $V\oplus\Der(\n)$ with bracket
\[
[(a,D),(b,E)]=(\{a,b\}+Db-Ea,\;DE-ED).
\]
Let $\g=(V,[\,,])$ be a Lie algebra on the same space and
$L\colon\g\to\Der(\n)$ a representation. Call a linear map
$P\colon V\to V$ a \emph{cocycle} (for $L$) if
\begin{equation}\label{eq:cocycle}
P[a,b]=\{Pa,Pb\}+L(a)Pb-L(b)Pa\qquad(a,b\in V),
\end{equation}
that is, if $x\mapsto(Px,L(x))$ is a Lie algebra homomorphism
$\g\to\h$. For a post-Lie algebra structure, \eqref{eq:pl1} says exactly
that $P=\mathrm{id}_V$ is a cocycle for $L$; together with \eqref{eq:pl2}
and \eqref{eq:pl3} this is \cite[Proposition 2.13]{BDV}.

Two trace identities will be used. First, \eqref{eq:pl1} gives
$R(x)=L(x)-\ad_\g(x)+\ad_\n(x)$. If $\n$ is nilpotent, each $\ad_\n(x)$
is nilpotent and has trace zero, so
\begin{equation}\label{eq:trace}
\tr R(x)=0\text{ for all }x\quad\Longleftrightarrow\quad
\tr L(x)=\tr\ad_\g(x)\text{ for all }x.
\end{equation}
Second, if $\g$ is perfect, every element of $\g$ is a sum of brackets,
and $L$ and $\ad_\g$ are representations, so $\tr L(x)=\tr\ad_\g(x)=0$
for all $x$.

Now let $\n$ be nilpotent and $z\in V$. The derivation $A=\ad_\h(z,0)$
of $\h$ is $A(a,D)=(\{z,a\}-Dz,0)$. It maps $\h$ into $V\oplus0$ and acts
there as $\ad_\n(z)$, so it is nilpotent, and $\exp(-A)$ is an
automorphism of $\h$ (a finite sum since $A$ is nilpotent, with
coefficients $1/j!$ since $k$ has characteristic zero).
Because $A$ has no $\Der(\n)$-component, for a cocycle $P$ the
homomorphism $\exp(-A)\circ(P,L)$ again has the form $x\mapsto(P^zx,L(x))$,
and $P^z$ is a cocycle. We call $P^z$ the \emph{conjugate} of $P$ by $z$.
For $\n$ abelian, $A^2=0$ and $P^z=P+\rho(z)$, where $\rho(z)x=L(x)z$;
in particular $\mathrm{id}^z=I+R(z)$.

\section{The determinant is invariant}

\begin{lemma}\label{lem:det}
Let $\n$ be nilpotent, let $\tr L(x)=\tr\ad_\g(x)$ for all $x\in V$, and
let $P$ be a cocycle with $\det P\ne0$. Then $\det P^z=\det P$ for every
$z\in V$.
\end{lemma}

\begin{proof}
Extend all maps $k[t]$-linearly to $V[t]=V\otimes k[t]$. The derivation
$tA$ of $\h[t]$ is nilpotent, and $\psi_t=\exp(-tA)\circ(P,L)$ is a
homomorphism $\g[t]\to\h[t]$ of the form $\psi_t(x)=(P_tx,L(x))$, with
$P_t$ a matrix over $k[t]$, $P_0=P$ and $P_1=P^z$. Differentiating in $t$
gives $\psi_t'=-A\psi_t$, whose first component is
\begin{equation}\label{eq:flow}
P_t'=-\ad_\n(z)\,P_t+\rho(z),\qquad \rho(w)\colon x\mapsto L(x)w.
\end{equation}
Let $F=\det P_t\in k[t]$ and $Q=\adj(P_t)$, so $QP_t=P_tQ=F\,I$. Since
$\psi_t$ is a homomorphism, $P_t$ satisfies \eqref{eq:cocycle} over
$k[t]$. Read as an identity of maps in $b$, this says
\[
P_t\ad_\g(a)=\bigl(\ad_\n(P_ta)+L(a)\bigr)P_t-\rho(P_ta)\qquad(a\in V[t]).
\]
Multiply on the right by $Q$ and take traces. Since
$\tr(P_t\ad_\g(a)Q)=\tr(\ad_\g(a)QP_t)=F\tr\ad_\g(a)$ and
$\tr\ad_\n(P_ta)=0$,
\[
\tr\bigl(\rho(P_ta)Q\bigr)=F\bigl(\tr L(a)-\tr\ad_\g(a)\bigr)=0
\qquad(a\in V[t]).
\]
For $w\in V[t]$ take $a=Qw$; then $P_ta=Fw$, so $F\tr(\rho(w)Q)=0$. As
$F(0)=\det P\ne0$, $F$ is a nonzero element of the domain $k[t]$, hence
$\tr(\rho(w)Q)=0$ for all $w$. Jacobi's formula $F'=\tr(QP_t')$ and
\eqref{eq:flow} now give
\[
F'=-\tr\bigl(Q\ad_\n(z)P_t\bigr)+\tr\bigl(Q\rho(z)\bigr)
=-F\tr\ad_\n(z)+0=0.
\]
In characteristic zero $F$ is constant, so
$\det P^z=F(1)=F(0)=\det P$.
\end{proof}

\section{Pushing a Levi factor down the lower central series}

Let $\n=\n^1\supseteq\n^2\supseteq\dots\supseteq\n^{r+1}=0$ be the lower
central series, $\n^{i+1}=\{\n,\n^i\}$, and recall that
$\{\n^i,\n^j\}\subseteq\n^{i+j}$ and that derivations preserve each
$\n^i$.

\begin{lemma}\label{lem:whitehead}
Let $\n$ be nilpotent, $\s\subseteq\g$ a semisimple subalgebra, and $P$ a
cocycle with $P(\s)\subseteq\n^i$ for some $i\ge1$. Then there is
$z\in\n^i$ with $P^z(\s)\subseteq\n^{i+1}$.
\end{lemma}

\begin{proof}
The space $M=\n^i/\n^{i+1}$ is an $\s$-module through $L$. For
$x,y\in\s$ we have $\{Px,Py\}\in\n^{2i}\subseteq\n^{i+1}$, so
\eqref{eq:cocycle} shows that $c\colon\s\to M$, $c(x)=Px+\n^{i+1}$,
satisfies $c([x,y])=L(x)c(y)-L(y)c(x)$. By Whitehead's first lemma there is
$w\in\n^i$ with $Px\equiv L(x)w\pmod{\n^{i+1}}$ for all $x\in\s$. Put
$z=-w$ and $A=\ad_\h(z,0)$. For $x\in\s$,
$A(Px,L(x))=(\{z,Px\}-L(x)z,0)$ has first component in $\n^i$, and each
further application of $A$ brackets with $z\in\n^i$, so $A^j(Px,L(x))$
has first component in $\n^{2i}\subseteq\n^{i+1}$ for $j\ge2$. Since
$\{z,Px\}\in\n^{i+1}$ too,
\[
P^zx\equiv Px+L(x)z=Px-L(x)w\equiv0\pmod{\n^{i+1}}.\qedhere
\]
\end{proof}

\section{Proofs of the main results}

\begin{proof}[Proof of Theorem~\ref{thm:main}]
Suppose $\g$ is not solvable. By Levi's theorem $\g$ has a nonzero
semisimple subalgebra $\s$. By \eqref{eq:trace},
$\tr L(x)=\tr\ad_\g(x)$ for all $x$. Start from the cocycle
$P^{(0)}=\mathrm{id}_V$, with $P^{(0)}(\s)\subseteq\n^1$. Applying
Lemma~\ref{lem:whitehead} $r$ times gives $z_1,\dots,z_r$ and cocycles
$P^{(i)}=(P^{(i-1)})^{z_i}$ with $P^{(i)}(\s)\subseteq\n^{i+1}$. By
Lemma~\ref{lem:det} and induction, $\det P^{(i)}=\det P^{(0)}=1$ for
every $i$. But $P^{(r)}(\s)\subseteq\n^{r+1}=0$ and $\s\ne0$, so
$P^{(r)}$ is singular, a contradiction.
\end{proof}

\begin{proof}[Proof of Corollary~\ref{cor:perfect}]
Suppose $(\g,\n)$ has a post-Lie algebra structure. As $\g$ is perfect,
$\tr L=\tr\ad_\g=0$ (Section~\ref{sec:setup}), so $\tr R=0$ by
\eqref{eq:trace}, and $\g$ is solvable by Theorem~\ref{thm:main}. A
solvable perfect Lie algebra is zero, since its derived series is
constant.
\end{proof}

\section{Formal verification}\label{sec:lean}
The theorem and corollary are proved in Lean~4 \cite{Lean} with Mathlib
\cite{Mathlib}. The definition of a post-Lie algebra structure, in
\texttt{PostLie/Definitions.lean}, takes two Lie algebras \texttt{g} and
\texttt{n} over a field \texttt{k} with \texttt{[CharZero k]}, both
finite-dimensional, and a linear equivalence \texttt{e : g ≃ₗ[k] n},
which plays the role of the common space $V$:
\begin{verbatim}
structure Product (e : g ≃ₗ[k] n) where
  mul : g →ₗ[k] g →ₗ[k] g
  skew : ∀ x y, mul x y - mul y x = ⁅x, y⁆ - e.symm ⁅e x, e y⁆
  left_lie : ∀ x y z, mul ⁅x, y⁆ z = mul x (mul y z) - mul y (mul x z)
  derivation : ∀ x y z, mul x (e.symm ⁅e y, e z⁆) =
    e.symm ⁅e (mul x y), e z⁆ + e.symm ⁅e y, e (mul x z)⁆
\end{verbatim}
Here \texttt{mul x y} is $x\cdot y$ and \texttt{e.symm ⁅e y, e z⁆} is
$\{y,z\}$, so the three fields are \eqref{eq:pl1}--\eqref{eq:pl3}. With
\texttt{right p x}, the map $y\mapsto y\cdot x$, the two results in
\texttt{PostLie/Statement.lean} are
\begin{verbatim}
theorem solvable_of_trace_right_eq_zero (e : g ≃ₗ[k] n) (p : Product e)
    [LieRing.IsNilpotent n]
    (htrace : ∀ x, LinearMap.trace k g (right p x) = 0) :
    LieAlgebra.IsSolvable g

theorem no_product_of_perfect (e : g ≃ₗ[k] n)
    [LieRing.IsNilpotent n] [Nontrivial g]
    (hperfect : ⁅(⊤ : LieIdeal k g), (⊤ : LieIdeal k g)⁆ = ⊤) :
    ¬ Nonempty (Product e)
\end{verbatim}
They are Theorem~\ref{thm:main} and Corollary~\ref{cor:perfect}:
\texttt{Nontrivial g} says $\g\ne0$, and \texttt{hperfect} says
$[\g,\g]=\g$. A build check prints the axioms each theorem depends on and
fails unless they are exactly Lean's \texttt{propext},
\texttt{Classical.choice} and \texttt{Quot.sound} and the following two.
\begin{itemize}
\item \texttt{PostLie.levi\_nonzero\_semisimple}: if \texttt{g} is not
solvable, it has a nontrivial subalgebra \texttt{s} with
\texttt{LieAlgebra.IsSemisimple k s}. This follows from Levi's theorem
\cite[Theorem 7.8.13]{We}: ``If $\g$ is a finite-dimensional Lie algebra
over a field of characteristic zero, then there is a semisimple Lie
subalgebra $\s$ of $\g$ (called a Levi factor of $\g$) such that $\g$ is
isomorphic to the semidirect product $\g=(\mathrm{rad}\,\g)\rtimes\s$.'' If
$\g$ is not solvable, then $\mathrm{rad}\,\g\ne\g$ and $\s\ne0$.
\item \texttt{PostLie.whitehead\_first}: for a finite-dimensional
semisimple \texttt{s}, a representation \texttt{ρ} of \texttt{s} on a
finite-dimensional \texttt{M} and a linear \texttt{c : s →ₗ[k] M} with
\texttt{c ⁅x, y⁆ = ρ x (c y) - ρ y (c x)}, there is \texttt{m} with
\texttt{c x = ρ x m} for all \texttt{x}. This is Whitehead's first lemma
\cite[Corollary 7.8.10]{We}: ``Let $\g$ be a semisimple Lie algebra over a
field of characteristic 0. If $M$ is any finite-dimensional $\g$-module,
then $H^1(\g,M)=0$. That is, every derivation from $\g$ into $M$ is an
inner derivation.''
\end{itemize}
Weibel calls a Lie algebra semisimple when its radical is zero, and these
are the finite products of simple Lie algebras \cite[Theorem 7.8.5]{We}.
Such a product satisfies Mathlib's definition, which asks that the ideal
lattice be generated by independent non-abelian atoms: the atoms are the
simple factors. Conversely, Mathlib proves that its semisimple Lie
algebras have zero radical
(\texttt{LieAlgebra.IsSemisimple.instHasTrivialRadical}). So each axiom is
an instance of the statement it quotes.
Everything else is proved, including Jacobi's formula over $k[t]$, the
exponential of the nilpotent derivation $A$ and the two lemmas. In the
formal proof a cocycle is a linear map from \texttt{g} to \texttt{n}, the
conjugation by $tz$ is the exponential series of $tA$ truncated to a
polynomial in $t$, whose value at $t=1$ is Mathlib's exponential of $A$,
and Lemma~\ref{lem:det} is the theorem \texttt{det\_translated}. To
reproduce the check, run \texttt{lake exe cache get} and \texttt{lake
build} in the \texttt{lean} folder of the accompanying files; with the
Mathlib cache the build takes about two minutes and under 3~GB of
memory.

\begin{thebibliography}{99}
\bibitem{Au} L. Auslander, Simply transitive groups of affine motions,
\emph{Amer. J. Math.} 99 (1977), 809--826.

\bibitem{Ba} O. Baues, Prehomogeneous affine representations and flat
pseudo-Riemannian manifolds, in \emph{Handbook of pseudo-Riemannian
geometry and supersymmetry}, IRMA Lect. Math. Theor. Phys. 16, European
Mathematical Society, 2010, 731--817,
\href{https://doi.org/10.4171/079-1/22}{doi:10.4171/079-1/22};
numbering of \href{https://arxiv.org/abs/0809.0824v1}{arXiv:0809.0824v1}.

\bibitem{Bu} D. Burde, Left-symmetric algebras, or pre-Lie algebras in
geometry and physics, \emph{Cent. Eur. J. Math.} 4 (2006), no.~3,
323--357,
\href{https://doi.org/10.2478/s11533-006-0014-9}{doi:10.2478/s11533-006-0014-9};
numbering and quotation from
\href{https://arxiv.org/abs/math-ph/0509016v2}{arXiv:math-ph/0509016v2}.

\bibitem{BD} D. Burde and K. Dekimpe, Post-Lie algebra structures on
pairs of Lie algebras, \emph{J. Algebra} 464 (2016), 226--245,
\href{https://doi.org/10.1016/j.jalgebra.2016.05.026}{doi:10.1016/j.jalgebra.2016.05.026};
numbering and quotations from
\href{https://arxiv.org/abs/1505.00955v1}{arXiv:1505.00955v1}.

\bibitem{BD13} D. Burde and K. Dekimpe, Post-Lie algebra structures and
generalized derivations of semisimple Lie algebras, \emph{Mosc. Math. J.}
13 (2013), no.~1, 1--18,
\href{https://doi.org/10.17323/1609-4514-2013-13-1-1-18}{doi:10.17323/1609-4514-2013-13-1-1-18};
numbering of \href{https://arxiv.org/abs/1108.5950v1}{arXiv:1108.5950v1}.

\bibitem{BDM} D. Burde, K. Dekimpe and M. Monadjem, Post-Lie algebra
structures for perfect Lie algebras, \emph{Comm. Algebra} 52 (2024),
no.~10, 4255--4267,
\href{https://doi.org/10.1080/00927872.2024.2344638}{doi:10.1080/00927872.2024.2344638};
numbering and quotations from
\href{https://arxiv.org/abs/2311.08985v1}{arXiv:2311.08985v1}.

\bibitem{BDV} D. Burde, K. Dekimpe and K. Vercammen, Affine actions on Lie
groups and post-Lie algebra structures, \emph{Linear Algebra Appl.} 437
(2012), no.~5, 1250--1263,
\href{https://doi.org/10.1016/j.laa.2012.04.007}{doi:10.1016/j.laa.2012.04.007};
numbering of \href{https://arxiv.org/abs/1109.0251v1}{arXiv:1109.0251v1}.

\bibitem{By} N. P. Byott, Solubility criteria for Hopf--Galois
structures, \emph{New York J. Math.} 21 (2015), 883--903.

\bibitem{Da} M. Damele, Determinant characters and multiplicative
solubility of finite skew braces, preprint, October 20, 2026,
\href{https://arxiv.org/abs/2610.28420v1}{arXiv:2610.28420v1}.

\bibitem{DL} M. Damele and A. Loi, Structural and rigidity properties of
Lie skew braces, \emph{J. Algebra} 695 (2026), 356--383,
\href{https://doi.org/10.1016/j.jalgebra.2026.01.048}{doi:10.1016/j.jalgebra.2026.01.048};
numbering and quotations from
\href{https://arxiv.org/abs/2507.06214v3}{arXiv:2507.06214v3}.

\bibitem{DO} J. Der\'e and M. Origlia, On post-Lie algebra structures
coming from simply transitive NIL-affine actions, \emph{Transform. Groups}
31 (2026), no.~2, 1391--1415,
\href{https://doi.org/10.1007/s00031-025-09906-2}{doi:10.1007/s00031-025-09906-2};
quotation from \href{https://arxiv.org/abs/2401.02503v1}{arXiv:2401.02503v1}.

\bibitem{Fo} D. J. F. Fox, Left-symmetric algebras and homogeneous
improper affine spheres, \emph{Ann. Global Anal. Geom.} 53 (2018), no.~3,
405--443,
\href{https://doi.org/10.1007/s10455-017-9582-0}{doi:10.1007/s10455-017-9582-0};
numbering and quotation from
\href{https://arxiv.org/abs/1707.08896v2}{arXiv:1707.08896v2}.

\bibitem{He} J. Helmstetter, Radical d'une alg\`ebre sym\'etrique \`a
gauche, \emph{Ann. Inst. Fourier} 29 (1979), no.~4, 17--35,
\href{https://doi.org/10.5802/aif.764}{doi:10.5802/aif.764}.

\bibitem{Kr} A. Kriebel, Post-Lie nonexistence with a semisimple target,
preprint, October 16, 2026,
\href{https://doi.org/10.5281/zenodo.23414025}{doi:10.5281/zenodo.23414025}.

\bibitem{Lean} L. de Moura and S. Ullrich, The Lean 4 theorem prover and
programming language, in \emph{Automated Deduction -- CADE 28}, Lecture
Notes in Comput. Sci. 12699, Springer, 2021, 625--635.

\bibitem{Mathlib} The mathlib Community, The Lean mathematical library, in
\emph{Proceedings of the 9th ACM SIGPLAN International Conference on
Certified Programs and Proofs}, 2020, 367--381.

\bibitem{OWR} D. Burde, Pre-Lie and post-Lie algebra structures (joint
work with K. Dekimpe), in \emph{Mini-Workshop: Poisson and Poisson-type
algebras}, \emph{Oberwolfach Rep.} 20 (2023), no.~4, 2688--2690,
\href{https://doi.org/10.4171/OWR/2023/46}{doi:10.4171/OWR/2023/46}.

\bibitem{Se} D. Segal, The structure of complete left-symmetric algebras,
\emph{Math. Ann.} 293 (1992), 569--578.

\bibitem{SV} A. Smoktunowicz and L. Vendramin, On skew braces (with an
appendix by N. Byott and L. Vendramin), \emph{J. Comb. Algebra} 2 (2018),
no.~1, 47--86,
\href{https://doi.org/10.4171/JCA/2-1-3}{doi:10.4171/JCA/2-1-3};
numbering and quotation from
\href{https://arxiv.org/abs/1705.06958v3}{arXiv:1705.06958v3}.

\bibitem{Va} B. Vallette, Homology of generalized partition posets,
\emph{J. Pure Appl. Algebra} 208 (2007), no.~2, 699--725,
\href{https://doi.org/10.1016/j.jpaa.2006.03.012}{doi:10.1016/j.jpaa.2006.03.012}.

\bibitem{We} C. A. Weibel, \emph{An Introduction to Homological Algebra},
Cambridge Stud. Adv. Math. 38, Cambridge University Press, 1994.
\end{thebibliography}
\end{document}
